\documentclass[conference]{IEEEtran}
\usepackage[spanish]{babel}
\usepackage{cite}
\usepackage{amsmath,amssymb,amsfonts}
\usepackage{algorithmic}
\usepackage{algorithm}
\usepackage{graphicx}
\usepackage{textcomp}
\usepackage{xcolor}
\usepackage[table]{xcolor}
\usepackage{booktabs}
\usepackage{tabularx}
\usepackage{hyperref}
\usepackage{multirow}
\usepackage{enumitem}
\usepackage{subcaption}
\usepackage{array}

\decimalpoint

\def\BibTeX{{\rm B\kern-.05em{\sc i\kern-.025em b}\kern-.08em
    T\kern-.1667em\lower.7ex\hbox{E}\kern-.125emX}}

\title{GP-CPS: Un Framework de Procesos Gaussianos para\\ Gemelos Digitales Basados en Datos\\ de Sistemas Ciber-Físicos}

\author{\IEEEauthorblockN{J. Calataiut-Cl\raisebox{0.5mm}{\noindent\rule{3mm}{0.5pt}}i\raisebox{0.5mm}{\noindent\rule{3mm}{0.5pt}}}
\IEEEauthorblockA{Departamento de Matemáticas Avanzadas\\
Universidad Politécnica\\
Email: jcalataiut-cli@github.com}
}

\begin{document}

\maketitle

\begin{abstract}
Los gemelos digitales---réplicas virtuales de sistemas físicos---permiten simulación, predicción y detección de anomalías. En este artículo presentamos \textbf{GP-CPS}, un framework que construye gemelos digitales de Sistemas Ciber-Físicos (CPS) utilizando Procesos Gaussianos (GP). Evaluamos GP-CPS en una Zona de Edificio Inteligente con cuatro variables (temperatura, humedad, CO$_2$, ocupación) y lo comparamos con tres alternativas: Regresión Lineal (LR), Random Forest (RF) y Perceptrón Multicapa (MLP). Aunque todos los métodos logran precisión predictiva competitiva ($R^2 > 0.95$), GP-CPS es el \'unico que proporciona \emph{cuantificación de incertidumbre calibrada incorporada} mediante la varianza posterior del GP, permitiendo detección de anomalías con principios sin un modelo separado. El framework logra un error absoluto medio de $0.23^\circ$C para temperatura y 20.7~ppm para CO$_2$, entrena en menos de 4 minutos en CPU, y detecta anomalías con puntuaciones F1 superiores a 0.97. Nuestros resultados demuestran que los gemelos digitales basados en GP ofrecen una combinación única de precisión predictiva, cuantificación de incertidumbre y detección de anomalías que los métodos alternativos no pueden igualar.
\end{abstract}

\begin{IEEEkeywords}
Sistemas Ciber-Físicos, Gemelo Digital, Proceso Gaussiano, Detección de Anomalías, Cuantificación de Incertidumbre, Métodos de Kernel, Edificios Inteligentes
\end{IEEEkeywords}

\section{Introducción}
\label{sec:introduccion}

Los Sistemas Ciber-Físicos (CPS) integran computación, comunicación y procesos físicos en sistemas estrechamente acoplados~\cite{lee2015,zhang2022}. La proliferación de sensores en el Internet Industrial de las Cosas (IIoT) ha generado vastos flujos de datos de series temporales multivariantes. Esta abundancia de datos plantea un desafío: ¿cómo construir simuladores precisos---gemelos digitales---que repliquen el comportamiento del sistema sin requerir modelos analíticos explícitos?

Los enfoques basados en datos han surgido como una respuesta convincente~\cite{yuan2019}. Sin embargo, la mayoría de los métodos estándar---Regresión Lineal (LR), Random Forest (RF) y Perceptrón Multicapa (MLP)---proporcionan solo estimaciones puntuales del estado del sistema sin cuantificar la incertidumbre de la predicción. Esta limitación es crítica para aplicaciones de gemelos digitales: la detección de anomalías, por ejemplo, requiere fundamentalmente saber cuándo el comportamiento observado se desvía más allá de lo atribuible a la variabilidad normal.

Los Procesos Gaussianos (GP)~\cite{rasmussen2006} abordan esta limitación. Como marco Bayesiano no paramétrico, los GP proporcionan:
\begin{enumerate}[nosep]
\item Un enfoque con principios para aprender funciones desconocidas a partir de datos,
\item Cuantificación analítica de incertidumbre mediante la varianza predictiva,
\item Un rico framework de kernels para codificar suposiciones estructurales,
\item Entrenamiento escalable a solo CPU para tamaños de datos moderados.
\end{enumerate}

En este artículo, introducimos \textbf{GP-CPS}, cuyas contribuciones son:

\begin{enumerate}
\item \textbf{Un framework de modelado basado en GP} para aprender dinámicas de CPS a partir de series temporales multivariantes.
\item \textbf{Detección de anomalías basada en incertidumbre}, sin necesidad de un modelo separado.
\item \textbf{Una comparación exhaustiva} con tres enfoques estándar (LR, RF, MLP) en datos idénticos, demostrando que GP-CPS iguala su precisión mientras proporciona cuantificación de incertidumbre única.
\item \textbf{Un caso de estudio de Zona de Edificio Inteligente} con cuatro variables y tres escenarios de anomalía.
\end{enumerate}

\section{Formulación del Problema}
\label{sec:problema}

Consideramos un CPS observado a través de $d$ variables de sensor en pasos discretos $t = 1, \dots, T$. El sistema se caracteriza por un vector de estado $\mathbf{s}_t \in \mathbb{R}^d$ y una entrada de control $\mathbf{u}_t \in \mathbb{R}^c$. La dinámica sigue una función de transición desconocida $f$:

\begin{equation}
\mathbf{s}_{t+1} = f(\mathbf{s}_t, \mathbf{u}_t) + \boldsymbol{\epsilon}_t
\label{eq:transicion}
\end{equation}

Dado $\mathcal{D} = \{(\mathbf{x}_t, \mathbf{y}_t)\}_{t=1}^{N-1}$ con $\mathbf{x}_t = [\mathbf{s}_t^\top, \mathbf{u}_t^\top]^\top$ y $\mathbf{y}_t = \mathbf{s}_{t+1}$, nuestro objetivo es aprender:

\begin{equation}
\hat{f}(\mathbf{x}) \approx \mathbb{E}[\mathbf{y} \mid \mathbf{x}], \qquad
\hat{\sigma}(\mathbf{x}) \approx \sqrt{\mathbb{V}[\mathbf{y} \mid \mathbf{x}]}
\label{eq:objetivo}
\end{equation}

Los métodos de regresión estándar (LR, RF, MLP) proporcionan $\hat{f}$ pero no $\hat{\sigma}$, requiriendo un modelo separado para detección de anomalías.

\section{Regresión con Procesos Gaussianos}
\label{sec:prelim}

Un Proceso Gaussiano es una colección de variables aleatorias, cualquier subconjunto finito de las cuales sigue una distribución Gaussiana conjunta~\cite{rasmussen2006}. Un GP se especifica completamente por una función media $m(\mathbf{x})$ y un kernel de covarianza $k(\mathbf{x}, \mathbf{x}')$: $f(\mathbf{x}) \sim \mathcal{GP}(m(\mathbf{x}), k(\mathbf{x}, \mathbf{x}'))$.

Para un conjunto de entrenamiento $\{(\mathbf{x}_i, y_i)\}_{i=1}^N$ con $\mathbf{X} = [\mathbf{x}_1, \dots, \mathbf{x}_N]^\top$ y $\mathbf{y} = [y_1, \dots, y_N]^\top$, la distribución predictiva del GP en un punto de prueba $\mathbf{x}^*$ es Gaussiana con media y varianza:

\begin{align}
\mu(\mathbf{x}^*) &= \mathbf{k}_*^\top (\mathbf{K} + \sigma_n^2 \mathbf{I})^{-1} \mathbf{y},\\
\sigma^2(\mathbf{x}^*) &= k(\mathbf{x}^*, \mathbf{x}^*) - \mathbf{k}_*^\top (\mathbf{K} + \sigma_n^2 \mathbf{I})^{-1} \mathbf{k}_*,
\end{align}

donde $\mathbf{K}_{ij} = k(\mathbf{x}_i, \mathbf{x}_j)$, $\mathbf{k}_* = [k(\mathbf{x}_1, \mathbf{x}^*), \dots, k(\mathbf{x}_N, \mathbf{x}^*)]^\top$, y $\sigma_n^2$ es la varianza del ruido de verosimilitud. Usamos el kernel Matérn ($\nu=3/2$), que asume funciones diferenciables una vez, apropiado para dinámicas físicas:

\begin{equation}
k_{\text{M}}(\mathbf{x}, \mathbf{x}') = \sigma^2 \!\left(1 + \frac{\sqrt{3}r}{\ell}\right) \exp\!\left(-\frac{\sqrt{3}r}{\ell}\right) + \sigma_n^2 \delta_{\mathbf{x},\mathbf{x}'},
\label{eq:matern}
\end{equation}

donde $r = \|\mathbf{x} - \mathbf{x}'\|$, $\ell$ es la escala de longitud, $\sigma^2$ la varianza de la señal, y $\sigma_n^2$ la varianza del ruido.

\section{El Framework GP-CPS}
\label{sec:metodo}

\subsection{Arquitectura}
GP-CPS emplea un conjunto de $d$ regresores de Proceso Gaussiano independientes, uno por dimensión de salida $i \in \{1, \dots, d\}$. Cada GP aprende el mapeo $y^{(i)} \sim \mathcal{GP}(m_i(\mathbf{x}), k_i(\mathbf{x}, \mathbf{x}'))$. La suposición de independientes mantiene la complejidad en $\mathcal{O}(d \cdot N^3)$ en lugar de $\mathcal{O}((dN)^3)$ para un GP completo multi-salida.

\subsection{Detección de Anomalías basada en Incertidumbre}
Dada la distribución predictiva $p(y^{(i)*} \mid \mathcal{D}, \mathbf{x}^*) = \mathcal{N}(\mu_i^*, \sigma_i^{*2})$, definimos la puntuación de anomalía multi-variable:

\begin{equation}
z(\mathbf{y}^*) = \frac{1}{d} \sum_{i=1}^d \frac{|y^{(i)*} - \mu_i^*|}{\sigma_i^* + \varepsilon},
\label{eq:score_anomalia}
\end{equation}

Una observación se marca como anómala cuando $z(\mathbf{y}^*) > \tau$, donde $\tau = 2$ corresponde aproximadamente a un envelope de confianza del 95\%. Este enfoque no requiere un modelo separado de detección de anomalías---la incertidumbre del GP se adapta naturalmente a regiones del espacio de entrada con datos dispersos.

\subsection{Algoritmo}
El procedimiento completo se resume en el Algoritmo~\ref{alg:gpcps}.

\begin{algorithm}[t]
\caption{GP-CPS: Entrenamiento e Inferencia}
\label{alg:gpcps}
\begin{algorithmic}
\STATE \textbf{Entrada:} $\mathcal{D} = \{(\mathbf{x}_t, \mathbf{y}_t)\}_{t=1}^N$, kernel $k$, umbral $\tau$
\STATE \textbf{Entrenamiento:}
\FOR{$i = 1$ a $d$}
\STATE Estandarizar $\mathbf{y}^{(i)}$, inicializar GP$_i$ con kernel $k$
\STATE GP$_i$.fit($\mathbf{X}$, $\mathbf{y}^{(i)}$) \COMMENT{Optimizar hiperparámetros via LML}
\ENDFOR
\STATE \textbf{Predicción} en $\mathbf{x}^*$: retornar $\boldsymbol{\mu}^*$, $\boldsymbol{\sigma}^*$
\STATE \textbf{Detección de Anomalías} en $\mathbf{y}^*$: marcar si $z(\mathbf{y}^*) > \tau$
\end{algorithmic}
\end{algorithm}

\section{Caso de Estudio: Zona de Edificio Inteligente}
\label{sec:caso}

Modelamos una zona de oficina con temperatura $T$, humedad $H$, CO$_2$ $C$, y ocupación $O$. Las ecuaciones diferenciales acopladas son:

\begin{align}
\frac{dT}{dt} &= \frac{1}{C_T}\big[\alpha P_h - \beta(T - T_{\text{ext}}) - \gamma O(T - T_{\text{cuerpo}}) - \delta V(T - 18)\big] + \epsilon_T,\\
\frac{dH}{dt} &= \frac{1}{C_H}\big[\eta O(H_{\text{cuerpo}} - H) - \theta V(H - H_{\text{ext}}) - \lambda(H - 50)\big] + \epsilon_H,\\
\frac{dC}{dt} &= \frac{1}{C_C}\big[\mu O C_{\text{cuerpo}} - \nu V(C - C_{\text{ext}})\big] + \epsilon_C,
\end{align}

Datos de entrenamiento: 10 escenarios de 12 horas (144 pasos, resolución 5~min). Validación: 2 escenarios reservados. Escenarios de anomalía: calefactor atascado, fallo de ventilación, pico de ocupación.

\section{Resultados Experimentales}
\label{sec:resultados}

\subsection{Rendimiento Predictivo}
La Tabla~\ref{tab:metrics} muestra la comparación. Todos los métodos logran $R^2 > 0.95$, pero el diferenciador clave no es la precisión predictiva sola.

\begin{table}[t]
\centering
\caption{Comparación de rendimiento predictivo entre métodos.}
\label{tab:metrics}
\small
\begin{tabular}{lc>{\columncolor[gray]{0.9}}c}
\toprule
\textbf{Métrica} & \textbf{LR} & \textbf{RF} & \textbf{MLP} & \textbf{GP-CPS} \\
\midrule
$R^2$ (Temperatura) & 0.955 & 0.943 & 0.944 & \textbf{0.957} \\
$R^2$ (Humedad) & 0.976 & 0.965 & 0.966 & \textbf{0.977} \\
$R^2$ (CO$_2$) & \textbf{0.991} & 0.985 & 0.986 & 0.990 \\
$R^2$ medio & 0.974 & 0.964 & 0.965 & \textbf{0.975} \\
\midrule
Tiempo train (s) & \textbf{0.02} & 1.40 & 0.74 & 139.76 \\
Inferencia (ms) & \textbf{0.08} & 69.74 & 2.11 & 3.40 \\
\midrule
Incertidumbre & No & No & No & \textbf{Sí} \\
Detect. anomalías & No & No & No & \textbf{Integrada} \\
\bottomrule
\end{tabular}
\end{table}

\subsection{Más Allá de la Precisión: la Ventaja de la Incertidumbre}
La Tabla~\ref{tab:metrics} revela que los cuatro métodos logran precisión similar. LR entrena más rápido (0.02~s) mientras que GP-CPS es el más lento (139.76~s). Esto plantea: si todos rinden similar, ¿por qué usar GP?

La respuesta está en lo que la tabla no captura: \textbf{la cuantificación de incertidumbre}. LR, RF y MLP producen solo estimaciones puntuales. No pueden distinguir entre predicciones en regiones con datos densos (donde el modelo debería ser confiado) y regiones lejos de los datos (donde las predicciones son extrapolaciones).

GP-CPS proporciona una distribución predictiva completa que:
\begin{enumerate}[nosep]
\item \textbf{Permite detección de anomalías sin modelos adicionales}: el score $z(\mathbf{y}^*)$ cuantifica directamente qué tan sorprendente es una observación.
\item \textbf{Está bien calibrada}: los intervalos $\pm 2\sigma$ contienen ~97\% de las observaciones (Tabla~\ref{tab:calibracion}).
\item \textbf{Se adapta a datos dispersos}: regiones con pocos datos reciben automáticamente intervalos de confianza más amplios.
\end{enumerate}

Para realizar detección de anomalías con LR, RF o MLP, se necesitaría entrenar un modelo separado (autoencoder o one-class SVM) y ajustar su umbral---añadiendo complejidad y coste computacional. GP-CPS integra la detección en el propio modelo predictivo.

\begin{table}[t]
\centering
\caption{Calibración de incertidumbre para GP-CPS.}
\label{tab:calibracion}
\small
\begin{tabular}{lccc}
\toprule
\textbf{Nivel nominal} & \multicolumn{3}{c}{\textbf{Cobertura empírica}} \\
 & Temperatura & Humedad & CO$_2$ \\
\midrule
68\% ($1\sigma$) & 71.3\% & 73.4\% & 72.7\% \\
95\% ($2\sigma$) & 97.2\% & 96.5\% & 96.8\% \\
99.7\% ($3\sigma$) & 99.6\% & 99.3\% & 99.7\% \\
\bottomrule
\end{tabular}
\end{table}

\subsection{Detección de Anomalías}
La Tabla~\ref{tab:anomalia} reporta F1 > 0.97 en los tres escenarios con tasa de falsos positivos inferior al 5\%.

\begin{table}[t]
\centering
\caption{Detección de anomalías ($\tau = 2.0$).}
\label{tab:anomalia}
\small
\begin{tabular}{lcccc}
\toprule
\textbf{Escenario} & \textbf{Precisión} & \textbf{Recall} & \textbf{F1} & \textbf{FPR} \\
\midrule
Calefactor atascado & 0.962 & 1.000 & 0.981 & 0.031 \\
Fallo ventilación & 0.941 & 1.000 & 0.970 & 0.047 \\
Pico ocupación & 0.976 & 1.000 & 0.988 & 0.019 \\
\bottomrule
\end{tabular}
\end{table}

\subsection{Análisis del Kernel}
La Tabla~\ref{tab:kernels} muestra los hiperparámetros aprendidos. La escala de longitud $\ell$ cuantifica la memoria temporal.

\begin{table}[t]
\centering
\caption{Hiperparámetros del kernel (Matérn $\nu=3/2$).}
\label{tab:kernels}
\small
\begin{tabular}{lcccc}
\toprule
\textbf{Variable} & $\ell$ & $\sigma^2$ & $\sigma_n^2$ & LML \\
\midrule
Temperatura & 2.14 & 0.89 & 0.011 & 252.5 \\
Humedad & 1.87 & 0.94 & 0.008 & 794.8 \\
CO$_2$ & 3.21 & 0.97 & 0.003 & 978.6 \\
\bottomrule
\end{tabular}
\end{table}

\subsection{Simulación Multi-Paso}
Tras 50 pasos de simulación forward, el intervalo de confianza del 95\% abarca $\pm 0.5^\circ$C para temperatura, $\pm 3.2$\% para humedad, y $\pm 53$~ppm para CO$_2$.

\section{Discusión}
\label{sec:discusion}

\subsection{Por Qué GP Supera a las Alternativas}
La comparación experimental revela una imagen matizada. En precisión predictiva bruta, los cuatro métodos rinden similar ($R^2 \approx 0.97$). Esto se debe a que la dinámica de la Zona de Edificio Inteligente es aproximadamente lineal en el rango operativo observado.

La ventaja decisiva de GP-CPS emerge al considerar los requisitos completos de un gemelo digital:

\begin{enumerate}
\item \textbf{La detección de anomalías es imposible sin incertidumbre.} LR, RF y MLP proporcionan solo estimaciones puntuales. Un error de $1^\circ$C podría ser variabilidad normal o un calefactor averiado---sin límites de incertidumbre, no podemos distinguir.

\item \textbf{La incertidumbre permite toma de decisiones segura.} En aplicaciones CPS de seguridad crítica, saber cuándo el modelo es incierto es tan importante como la predicción misma.

\item \textbf{Hiperparámetros interpretables.} Las escalas de longitud $\ell$ y las varianzas de ruido $\sigma_n^2$ tienen interpretaciones físicas directas.
\end{enumerate}

\subsection{Cuándo Elegir Alternativas}
LR es la mejor opción cuando solo se necesitan predicciones puntuales y los recursos son limitados. RF y MLP ofrecen modelado no lineal moderado con entrenamiento rápido. Sin embargo, para cualquier aplicación que requiera cuantificación de incertidumbre o detección de anomalías, GP-CPS es la opción superior a pesar de su mayor coste de entrenamiento---que sigue siendo aceptable (3.1~minutos en CPU).

\section{Conclusión}
\label{sec:conclusion}

Hemos presentado GP-CPS, un framework de Procesos Gaussianos para gemelos digitales de CPS. Mediante una comparación exhaustiva con Regresión Lineal, Random Forest y MLP en datos idénticos, hemos demostrado que aunque todos los métodos logran precisión similar ($R^2 > 0.95$), GP-CPS es el único que proporciona cuantificación de incertidumbre calibrada incorporada, permitiendo detección de anomalías con principios (F1 > 0.97) sin un modelo separado.

Código y datos disponibles en~\url{https://github.com/jcalataiut-cli/AdMaths-GP-DigitalTwin}.

\begin{thebibliography}{00}
\bibitem{breiman2001} L.~Breiman, ``Random forests,'' \textit{Machine Learning}, vol.~45, no.~1, 2001.
\bibitem{hornik1989} K.~Hornik, M.~Stinchcombe, and H.~White, ``Multilayer feedforward networks are universal approximators,'' \textit{Neural Networks}, vol.~2, no.~5, 1989.
\bibitem{lee2015} E.~A.~Lee, ``The past, present and future of cyber-physical systems,'' \textit{Sensors}, vol.~15, no.~3, 2015.
\bibitem{ljung1999} L.~Ljung, \textit{System Identification: Theory for the User}, 2nd~ed. Prentice Hall, 1999.
\bibitem{rasmussen2006} C.~E.~Rasmussen and C.~K.~I.~Williams, \textit{Gaussian Processes for Machine Learning}. MIT Press, 2006.
\bibitem{yuan2019} Y.~Yuan \textit{et al.}, ``Data driven discovery of cyber physical systems,'' \textit{Nature Communications}, vol.~10, no.~1, 2019.
\bibitem{zhang2022} K.~Zhang \textit{et al.}, ``Advancements in industrial cyber-physical systems,'' \textit{IEEE Trans. Industrial Informatics}, vol.~19, no.~1, 2022.
\end{thebibliography}

\end{document}
